\chapter{Hydromechanical behaviour of biochar-amended soil along wetting and drying paths}
\section{Introduction}
The term water retention behaviour in this chapter is used to describe the ability of biochar-amended soil to retain, release, and transmit water under unsaturated conditions. This behaviour is quantified through the soil-water retention relationship and the unsaturated hydraulic conductivity function. The former defines the variation of water content or degree of saturation with suction, while the latter describes the conductivity of the partially saturated pore system as a function of suction or water content. In addition, the water retention behaviour considered here is not treated as an intrinsic material response independent of soil fabric; rather, it is interpreted as a pore-structure- and surface-property-dependent response that may vary with drying–wetting history, initial compaction degree, and applied stress state. Mercury intrusion porosimetry and contact-angle measurements are therefore used to support the interpretation of the measured retention and conductivity responses. \par% the basic definition should be mentioned in the beginning of the introduction. please check! or this paragraph should be mentioned in the introdution section.

it should mentioned that in the introdution or inthe summary or conclusion, the aging effect of biochar on the soil hydromechancial behaviours is not considered in this chapter. \par % A relative literatures is: Influence of biochar on the soil water retention characteristics (SWRC): Potential application in geotechnical engineering structures

It should be noted that the sign convention adopted in this chapter takes compressive pressure as positive.\par

The following problems are studied in this chapter:
\begin{itemize}
    \item How does the addition of biochar affect the soil water retention characteristics (SWRC)
    \item Hysteresis behaviour of SWRC with wetting and drying cycles
    \item Hydraomechnical behaviours under different wetting methods
\end{itemize} 

\section{Materials}
\subsection{Soil}
The tested soil was collected near the UK \ac{HS2} construction site. Bulk soils were excavated within a depth of 1.0 meter below the ground surface. Before testing, visible gravels and plant roots were carefully removed manually. The soil was then oven-dried at \SI{105}{\degreeCelsius} for 48 hours and subsequently passed through a 2 mm sieve for testing. The basic properties were determined following the BS EN ISO 17892 and BS 1377 standard series. Figure \ref{fig:ch3GrainSizeDistribution}a presents its grain size distribution. A summary of tested soil's properties is provided in Table \ref{tab:ch3soilproperties}. According to the BS EN ISO 14688-2, the soil exhibits a relative broad grain size distribution. It is classified as a low plasticity SAND with silt and clay (CIL). \par
\begin{figure}[htbp]
    \centering
    \includegraphics[width=1.0\linewidth]{../figures/Ch3/ch3GrainSize.png}
    \caption{Grain size distribution of tested materials: (a) soil and (b) biochar}
    \label{fig:ch3GrainSizeDistribution}
\end{figure}

\begin{table}[htbp]
    \centering
    \caption{Basic geotechnical properties of the tested soil}
    \label{tab:ch3soilproperties}
    \begin{tabularx}{\textwidth}{>{\raggedright\arraybackslash}X >{\raggedright\arraybackslash}p{0.40\textwidth}}
        \toprule
        Parameter & Value \\
        \midrule
        Specific gravity, $G_{\mathrm{s,s}}$ & 2.6\\
        Organic matter content, $w_{\mathrm{org}}$ (\%) & 3.62 \\
        \multicolumn{2}{l}{\textit{Grading (BS EN ISO 17892-4 and BS EN ISO 14688-2)}} \\
        Sand content (\SI{63}{\micro\meter} to \SI{2}{\milli\meter}, \%) & 76.48 \\
        Silt content (\SI{2}{\micro\meter} to \SI{63}{\micro\meter}, \%) & 22.02 \\
        Clay content (\textless\SI{2}{\micro\meter}, \%) & 1.5 \\
        $D_{10}$ ($\mu$m) & 17.62 \\
        $D_{30}$ ($\mu$m) & 81.07 \\
        $D_{60}$ ($\mu$m) & 218.71 \\
        Uniformity coefficient, $C_U$ & 12.41 \\
        Coefficient of curvature, $C_C$ & 1.71 \\
        \multicolumn{2}{l}{\textit{Atterberg limits (BS 1377-2)}} \\
        Liquid limit, $w_L$ (\%) & 31.54 \\
        Plastic limit, $w_P$ (\%) & 18.79 \\
        Plasticity index, $I_P$ (\%) & 12.75 \\
        \multicolumn{2}{l}{\textit{Compaction test (BS 1377-2)}} \\
        Maximum dry density, $\rho_{\mathrm{d,max}}$ (g/cm$^3$) & 1.80 \\
        Optimum moisture content, $w_{\mathrm{opt}}$ (\%) & 13.9 \\
        \bottomrule
    \end{tabularx}
\end{table}

\subsection{Biochar}
The tested biochar, provided by a commercial supplier, is derived from arboricultural arisings comprising mixed European hardwood species, including oak and robinia. It is produced through pyrolysis for 4 to 8 h in a retort kiln at around \SI{425}{\degreeCelsius}. The tapped bulk density, measured using a standardised tapping procedure, is 380 kg/m$^3$. Its specific gravity ($G_{\mathrm{s,b}}$) is 1.07. Due to the low massive density and friable nature of the biochar, a Bettersize 2600 laser particle size analyser operating in wet mode is employed to determine its particle size distribution. It should be noted that this measured particle size distribution is expressed as a volume percentage. As presented in Figure \ref{fig:ch3GrainSizeDistribution}b, the particle size of biochar ranges from \SI{1.57}{\micro\meter} to \SI{1,826}{\micro\meter}. The incremental passing curve exhibits a significant bimodal pattern with two dominant particle sizes of \SI{98.09}{\micro\meter} and \SI{631.41}{\micro\meter}.\par

\subsection{Soil-biochar mixture}
There are three biochar dosages being considered, 0\%, 5\%, and 10\%, which are the relative mass percentages of biochar to the dry soil. The specific gravities of soil-biochar mixture are 2.6, 2.43, and 2.3 for 0\%, 5\%, and 10\% biochar dosages, respectively, which are estimated via:
\begin{equation}
    G_{\mathrm{s}} = \frac{(1 + \alpha)}{1/G_{\mathrm{s,s}} + \alpha/G_{\mathrm{s,b}}}
\end{equation}
where $\alpha$ is the biochar dosage. Standard compaction tests were conducted to determine the \ac{MDD} and \ac{OMC} at different biochar dosages. The test results are shown in Figure \ref{fig:ch3MDDOMC}. The \ac{OMC} increased with increasing biochar dosage, while the \ac{MDD} decreased. The \acp{MDD} correspond to the air voids of 18.4\%, 16.5\%, and 13\%, respectively.\par
\begin{figure}[htbp]
    \centering
    \includegraphics[width=1.0\linewidth]{../figures/Ch3/ch3MDDOMC.png}
    \caption{Maximum dry density and optimum moisture content of biochar-amended soil at different biochar dosages: (a) pure soil, (b) 5\% biochar, and (c) 10\% biochar}
    \label{fig:ch3MDDOMC}
\end{figure}

\section{Test programme and apparatus}
\subsection{Test programme}
Three series of main laboratory tests were conducted to investigate the hydromechanical behaviour of biochar-amended soils. The first series comprised evaporation tests designed to evaluate the influence of the initial compaction state (dense or loose) on the \ac{SWRC} and hydraulic conductivity during drying. The second series focused on the \ac{SWRC} and its hysteretic behaviour under different vertical net stresses, with particular attention to the effects of \ac{WDC}. The third series involved suction-controlled oedometer tests performed at different matric suctions to assess the mechanical response of the specimens. In addition, complementary physicochemical and microstructural characterisation techniques were employed to support the interpretation of the experimental results and to elucidate the underlying mechanisms governing the observations.\par
\subsection{Tensiometer and dew point potentiometer}
The first series of tests was conducted by using a combined tensiometer and dew-point hygrometry method. The test can be implemented using a HYPROP system for continuous measurements and a WP4C dew-point potentiometer for discrete measurements \parencite{singhEstimatingSoilWater2024, satyanagaMeasurementSoilwaterCharacteristic2019}. Both instruments were supplied by METER Group (USA).\par
As shown in Figure \ref{fig:ch3WP4CHyprop2}, the HYPROP system consists of tensiometer sensors, a specimen ring, an electronic balance, and a computer for data acquisition and analysis. The specimen is 80 mm in diameter and 50 mm in height. During the test, water loss from the specimen is continuously measured and automatically recorded by the electronic balance. Two vertically aligned tensiometers are installed within the specimen to continuously monitor matric suction, with their ceramic tips positioned 12.5 mm above and below the specimen centre, respectively. The maximum matric suction that can be reliably measured by the HYPROP system is approximately 100 kPa. To extend the measurement range, a WP4C device is employed to determine the higher suction values through the chilled-mirror dew-point technique \parencite{changWaterTransferGranite2026, shokranaMeasurementSoilWater2020}. It should be noted that the WP4C measures total suction, which comprises both matric and osmotic suction components. However, the osmotic suction is generally negligible for the soil tested in this study, and previous studies have demonstrated good agreement between WP4C and tensiometer measurements \parencite{aziziCouplingCyclicWater2023}. \par
The unsaturated hydraulic conductivity is derived using the simplified evaporation method \parencite{petersSimplifiedEvaporationMethod2008, schindlerSimplifyingEvaporationMethod2006}, in which the water flux ($q_i$) is calculated from the measured change in water volume as:
\begin{equation}
    q_i=\frac{\Delta V_i}{2A\Delta t_i}
\end{equation}
where $\Delta V_i$ is the water volume loss over the time interval $\Delta t_i$, and $A$ is the cross-sectional area of the specimen. The hydraulic conductivity ($k_{\mathrm{un,}i}$) can be obtained based on Darcy's law as:
\begin{equation}
    k_{\mathrm{un,}i}=-\frac{q_i}{(\Delta s_i/\Delta z)-1}
\end{equation}
where $\Delta s_i$ is the suction difference between the two tensiometer tips and $\Delta z$ is their vertical spacing.\par

\begin{figure}[htbp]
    \centering
    \includegraphics[width=0.88\linewidth]{../figures/Ch3/ch3WP4C&Hyprop2.png}
    \caption{HYPROP (left) and WP4C (right) apparatuses}
    \label{fig:ch3WP4CHyprop2}
\end{figure}

\subsection{Suction-controlled oedometer apparatus} \label{sec:ch3oedometerapparatus}
The other two series of tests were conducted on a suction-controlled oedometer apparatus (VJ Tech, UK). As shown in Figure \ref{fig:ch3apparatusSWRC}, the apparatus consists of a loading frame, a pneumatic air pressure controller, a pore water pressure controller, a flushing system, and a data acquisition system. The loading frame is equipped with a loading unit and a displacement transducer to apply a predefined vertical stress and measure the resulting vertical displacement, respectively. The sample is in a diameter of 70 mm and a height of 20 mm. The air-pressure and pore-water-pressure controllers, together with a pre-saturated \ac{HAEV} ceramic disc (AEV = 500 kPa), are used to impose target matric suctions on the specimen according to the axis-translation technique \parencite{hilfInvestigationPorewaterPressure1956}. Pore-air pressure and pore-water pressure are applied above and below the \ac{HAEV} ceramic disc, respectively. The vertical net pressure ($p_{\mathrm{net}}$) and the matrix suction ($s$) applied to the specimen is given by: 
\begin{equation}
    p_{\mathrm{net}} = p_\mathrm{v} - p_\mathrm{a}
\end{equation}
\begin{equation}
    s = p_\mathrm{a} - p_\mathrm{w}
\end{equation}
where $p_{\mathrm{v}}$ is the applied vertical pressure, $p_\mathrm{a}$ is the air pressure, and $p_\mathrm{w}$ is the pore water pressure. The flushing system is used to pre-saturate the \ac{HAEV} disc and to flush out dissolved air that exsolves beneath the \ac{HAEV} disc. The outflow water was carefully collected.\par
\begin{figure}[htbp]
    \centering
    \includegraphics[width=0.76\linewidth]{../figures/Ch3/ch3apparatusSWRC.png}
    \caption{Suction-controlled oedometer apparatus}
    \label{fig:ch3apparatusSWRC}
\end{figure}

\subsection{Surface and microstructural characterisation apparatus}
\Ac{XRD} analysis was conducted to detect mineral composition using a SmartLabXE X-ray diffractometer. The diffraction pattern was collected over a \(2\theta\) range of \ang{10}--\ang{90} with a step size of \ang{0.01}.\par
The surface hydrophobicity of tested materials was reflected via static and dynamic water contact angles, which were measured using the sessile drop method. The test was conducted on a Lauda Scientific LSA100S-T instrument (Germany). Before testing, the dried biochar and soil was ground into a fine powder and sieved through a 200-mesh screen. Approximately 200 mg of the powder was then compressed into a flat, thin pellet using a hydraulic press at 15 MPa for 1 minute. A 3 $\mu$L droplet of deionised water was carefully deposited onto the pellet surface using a precision microsyringe. Droplet images were recorded with a high-speed camera and analysed using the Young-Laplace model. To ensure measurement reliability, each sample was tested in triplicate, and the average contact angle values were reported. \par
\Ac{SEM} was performed using a Philips XL30 ESEM in secondary electron mode to examine the microstructure and interfacial interactions between soil and biochar. Before scaning, samples were gold-coated to enhance conductivity.\par
The \ac{PSD} of the specimens was determined by \ac{MIP} mothod using an AutoPore V9600 porosimeter (Micromeritics Inc., USA). Carefully trimmed specimens were first immersed in liquid nitrogen and then freeze-dried for 72 hours using a Labconco FreeZone freezing dryer to ensure complete drying. Freeze-drying was adopted because it is effective in preserving the original pore structure of the specimens \parencite{ngEffectsPoreStructure2023}. Owing to the non-wetting nature of mercury, its intrusion requires to overcome the interfacial resistance exerted by the pore surfaces of the soil and biochar. By idealising the pores within the specimen as cylindrical capillaries, the equilibrium relationship between the equivalent pore radius ($r_{\mathrm{eq}}$) and the required pressure ($p_{\mathrm{m}}$) is given by the Washburn equation as:
\begin{equation}
    r_{\mathrm{eq}} = -\frac{2\gamma \cos \theta}{p_{\mathrm{m}}}
\end{equation}
where $\gamma$ is the surface tension of mercury and $\theta$ is the contact angle between mercury and the pore surface. The volume of intrusing mercury at each pressure step is recorded, and the corresponding pore-size distribution can be derived. \par

\section{Sample preparation and test procedures} \label{sec:ch3testdetails}
\subsection{Sample preparation}
To prepare the samples, dry soil and biochar were first weighted and thoroughly mixed, after which de-aired water was added to achieve the target moisture content. The mixed materials were then sealed in plastic bags and stored for at least 24 hours to ensure moisture equilibration. All the tests were conducted on recompacted specimens. Specimens were prepared at their respective \acp{OMC} using the undercompaction method \parencite{laddPreparingTestSpecimens1978}. Two compaction states, namely dense and loose conditions, were considered, corresponding to \ac{DC} of 90\% and 80\%, respectively. Here, the \ac{DC} is defined as the ratio of the specimen dry density to the corresponding \ac{MDD}. In the first series of tests, both dense and loose specimens were tested. In the second and third series, only loose specimens were prepared. All tests were conducted with biochar dosages of 0\% (untreated soil), 5\%, and 10\%. The initial states of prepared specimens are summarised in Table \ref{tab:ch3_test_programme_states}.\par
Six parallel specimens, marked with asterisk in Table \ref{tab:ch3_test_programme_states}, were prepared for the \ac{MIP}. S1B0D and S1B5D were selected as representative examples for \ac{SEM} test to illustrated the microstructural characteristics of the soil matrix and biochar particles. These eight specimens were freeze-dried for before testing. \par
\afterpage{%
\clearpage
\begingroup
\begin{landscape}
\centering
\scriptsize
\setlength{\tabcolsep}{2pt}
\renewcommand{\arraystretch}{1.15}
\setlength{\LTleft}{0pt}
\setlength{\LTright}{0pt}
\newcommand{\dosagecell}[2]{\multirow{##1}{=}{\centering\arraybackslash ##2}}
\begin{longtable}{@{}p{36pt}p{62pt}>{\centering\arraybackslash}p{50pt}p{70pt}p{47pt}p{70pt}p{70pt}p{47pt}p{70pt}p{70pt}@{}}
    \caption{Summary of the specimen details}
    \label{tab:ch3_test_programme_states}\\
    \toprule
    \multirow{3}{=}{Test series} & \multirow{3}{=}{Test ID} & \multirow{3}{=}{\makecell[l]{Biochar\\dosage (\%)}} & \multirow{3}{=}{Loading path} & \multicolumn{3}{c}{After compaction} & \multicolumn{3}{c}{After saturation} \\
    \cmidrule(lr){5-7}\cmidrule(l){8-10}
                                 &                          &                                               &                               & \makecell[l]{Void ratio\\($e_0$)} & \makecell[l]{Gravimetric\\moisture\\content\\($w_0$, \%)} & \makecell[l]{Degree of\\saturation\\($S_{\mathrm{r},0}$, \%)} & \makecell[l]{Void ratio\\($e_{\mathrm{s}}$)} & \makecell[l]{Gravimetric\\moisture\\content\\($w_{\mathrm{s}}$, \%)} & \makecell[l]{Degree of\\saturation\\($S_{\mathrm{r},\mathrm{s}}$, \%)} \\
    \midrule
    \endfirsthead
    \caption[]{Summary of the specimen details (continued)}\\
    \toprule
    \multirow{3}{=}{Test series} & \multirow{3}{=}{Test ID} & \multirow{3}{=}{\makecell[l]{Biochar\\dosage (\%)}} & \multirow{3}{=}{Loading path} & \multicolumn{3}{c}{After compaction} & \multicolumn{3}{c}{After saturation} \\
    \cmidrule(lr){5-7}\cmidrule(l){8-10}
                                 &                          &                                               &                               & \makecell[l]{Void ratio\\($e_0$)} & \makecell[l]{Gravimetric\\moisture\\content\\($w_0$, \%)} & \makecell[l]{Degree of\\saturation\\($S_{\mathrm{r},0}$, \%)} & \makecell[l]{Void ratio\\($e_{\mathrm{s}}$)} & \makecell[l]{Gravimetric\\moisture\\content\\($w_{\mathrm{s}}$, \%)} & \makecell[l]{Degree of\\saturation\\($S_{\mathrm{r},\mathrm{s}}$, \%)} \\
    \midrule
    \endhead
    \midrule
    \multicolumn{10}{r}{Continued on next page}\\
    \endfoot
    \bottomrule
    \multicolumn{10}{@{}p{628pt}@{}}{Note: (1) The tests marked with an asterisk are also prepared for the \ac{MIP} test. The Test ID indicates the test series, biochar dosage, and stress path. (2) In the first series test, "D" and "L" denote dense and loose recompacted specimens, corresponding to degrees of compaction of 90\% and 80\%, respectively. (3) The void ratio in Test Series I is assumed to be constant during the saturation process, i.e., $e_{\mathrm{s}} = e_{0}$.}\\
    \endlastfoot
    \multirow{6}{=}{I}
        & S1B0D*   & \dosagecell{2}{0}    & \multirow{6}{=}{\makecell[l]{Free\\evaporation}} & 0.599 & 13.94 & 60.5 & 0.599 & 22.99 & 99.8 \\
        & S1B0L*   &                     &                                                                            & 0.802 & 14.65 & 47.5 & 0.802 & 30.85 & 100 \\
    \cmidrule(lr){3-3}
        & S1B5D*   & \dosagecell{2}{5}    &                                                                            & 0.59 & 15.64 & 64.4 & 0.59 & 24.09 & 99.2 \\
        & S1B5L*   &                     &                                                                            & 0.803 & 15.3 & 46.3 & 0.803 & 33.05 & 100 \\
    \cmidrule(lr){3-3}
        & S1B10D*  & \dosagecell{2}{10}   &                                                                            & 0.603 & 16.47 & 62.8 & 0.603 & 25.72 & 98.1 \\
        & S1B10L*  &                     &                                                                            & 0.807 & 16.33 & 46.6 & 0.807 & 35.04 & 100 \\
    \midrule
    \multirow{12}{=}{II}
        & S2B0V0    & \dosagecell{4}{0}    & HV0   & 0.785 & 13.58 & 45 & 0.783 & 30.12 & 100 \\
        & S2B0V50   &                     & HV50  & 0.784 & 14.59 & 48.4 & 0.784 & 30.12 & 99.9 \\
        & S2B0V100  &                     & HV100 & 0.784 & 14.5 & 48.1 & 0.78 & 30 & 100 \\
        & S2B0V200  &                     & HV200 & 0.791 & 13.37 & 43.9 & 0.787 & 30.27 & 100 \\
    \cmidrule(lr){3-3}
        & S2B5V0    & \dosagecell{4}{5}    & HV0   & 0.805 & 15.38 & 46.4 & 0.802 & 32.51 & 98.5 \\
        & S2B5V50   &                     & HV50  & 0.792 & 14.57 & 44.7 & 0.791 & 32.55 & 100 \\
        & S2B5V100  &                     & HV100 & 0.794 & 15.76 & 48.2 & 0.794 & 32.67 & 100 \\
        & S2B5V200  &                     & HV200 & 0.801 & 16.07 & 48.8 & 0.801 & 32.63 & 99 \\
    \cmidrule(lr){3-3}
        & S2B10V0   & \dosagecell{4}{10}   & HV0   & 0.795 & 16.9 & 48.9 & 0.795 & 34.12 & 98.7 \\
        & S2B10V50  &                     & HV50  & 0.787 & 16.37 & 47.8 & 0.786 & 34.17 & 100 \\
        & S2B10V100 &                     & HV100 & 0.797 & 17.2 & 49.6 & 0.792 & 34.43 & 100 \\
        & S2B10V200 &                     & HV200 & 0.794 & 17.56 & 50.9 & 0.794 & 33.93 & 98.3 \\
    \pagebreak
    \multirow{12}{=}{III}
        & S3B0S100   & \dosagecell{4}{0}    & CS100 & 0.796 & 14.37 & 46.9 & 0.796 & 30 & 98 \\
        & S3B0S200   &                     & CS200 & 0.792 & 14.53 & 47.7 & 0.791 & 30.42 & 100 \\
        & S3B0S300   &                     & CS300 & 0.799 & 13.98 & 45.5 & 0.795 & 30.33 & 99.2 \\
        & S3B0S400   &                     & CS400 & 0.784 & 14.27 & 47.3 & 0.783 & 30.12 & 100 \\
    \cmidrule(lr){3-3}
        & S3B5S100   & \dosagecell{4}{5}    & CS100 & 0.787 & 15.95 & 49.2 & 0.787 & 32.39 & 100 \\
        & S3B5S200   &                     & CS200 & 0.792 & 14.78 & 45.3 & 0.788 & 32.01 & 98.7 \\
        & S3B5S300   &                     & CS300 & 0.791 & 15.51 & 47.6 & 0.79 & 31.86 & 98 \\
        & S3B5S400   &                     & CS400 & 0.795 & 15.36 & 46.9 & 0.792 & 32.59 & 100 \\
    \cmidrule(lr){3-3}
        & S3B10S100  & \dosagecell{4}{10}   & CS100 & 0.79 & 17.18 & 50 & 0.79 & 34 & 99 \\
        & S3B10S200  &                     & CS200 & 0.795 & 17.01 & 49.2 & 0.795 & 34.36 & 99.4 \\
        & S3B10S300  &                     & CS300 & 0.795 & 17.31 & 50.1 & 0.795 & 34.57 & 100 \\
        & S3B10S400  &                     & CS400 & 0.786 & 17.3 & 50.6 & 0.786 & 34.17 & 100 \\
\end{longtable}
\end{landscape}
\clearpage
\endgroup
}

\subsection{Test procedures}
In the first series of tests, specimens were vacuum-saturated. The sensors and tensiometer shafts were also fully saturated prior to installation. During the test, the specimen was allowed to evaporate naturally under laboratory conditions at approximately 22 °C. The average matric suction was calculated as the mean value measured by the two tensiometers. After the measured matric suction approached the cavitation limit of the tensiometer shafts at around 100 kPa, the test was transferred to the WP4C device. At this stage, the specimen was removed from the HYPROP system and carefully cut into blocks of approximately 2 cm. These sample blocks were then subjected to continued evaporation under the same laboratory conditions. During evaporation, selected blocks were tested using the WP4C to determine the water potential, after which their water contents were measured immediately.\par
For the \ac{WDC} tests, the prepared specimen was first placed in the sample cell. The matric suction was set to 0 kPa to bring the specimen to a saturated state, while a very small vertical net stress was maintained to ensure close contact between the loading platen and the specimen surface, allowing axial deformation to be monitored. At this stage, the specimen corresponded to point O in Figure \ref{fig:ch3loadingpath}a. The matric suction was then maintained, and the vertical net stress was increased to the target values (points P2, P3, and P4, respectively) until equilibrium was reached. Subsequently, the specimen was dried stepwise to a matric suction of 400 kPa (0 kPa $\to$ 100 kPa $\to$ 200 kPa $\to$ 300 kPa $\to$ 400 kPa), and was then wetted along the same path. Four loading paths were conducted, corresponding to vertical net stresses of 0, 50, 100, and 200 kPa. Therefore, the hysteresis in the\ac{SD-SWRC} can be investigated. After completion of a \ac{WDC}, the specimen was subjected to stepwise vertical loading up to 200 kPa, reaching point P4. During the test, the vertical displacement and inflow/outflow water volumewere continuously monitored. Equilibrium at each loading step was considered to have been achieved when the daily inflow/outflow rate was lower than 0.04\% \parencite{ngWaterRetentionVolumetric2016a,muHydromechanicalBehaviorUnsaturated2023}.\par
The suction-controlled oedometer test started from the same loading state, i.e., saturated condition at point O in Figure \ref{fig:ch3loadingpath}b. Four loading paths were considered. The specimens were first dried to the target matric suctions, corresponding to points O1, O2, O3, and O4, respectively. After suction equilibrium was reached, vertical net stress was applied stepwise to 50, 100, 200, 400, and 800 kPa, leading to points X1, X2, X3, and X4. Subsequently, the matric suction was progressively reduced to 0 kPa while remaining a net pressure of 800 kPa, finally reaching point X. The variations in void ratio and moisture content were determined from the measured vertical displacement and the volume of outflowing or inflowing water. Equilibrium was considered when the rates of volumetric strain and water flow were lower than 0.025\% and 0.04\% per day, respectively. \parencite{ngEffectsSoilStructure2017, muHydromechanicalBehaviorUnsaturated2023}. \par
\begin{figure}[htbp]
    \centering
    \includegraphics[width=1.0\linewidth]{../figures/Ch3/ch3loadingpath.png}
    \caption{Loading path for suction-controlled oedometer test: (a) wetting and drying cycle test under constant vertical net pressure, and (b) suction-controlled oedometer test}
    \label{fig:ch3loadingpath}
\end{figure}

\section{Results and discussion}
\subsection{Contact angle} \label{sec:ch3contactangle}
Figure \ref{fig:ch3CA} presents the contact angle measurements of the tested materials. Pure biochar exhibited water repellency, whereas the untreated soil was strongly hydrophilic, with a contact angle of only 15.5°. The hydrophobic nature of biochar is commonly attributed to aromatic carbon structures and non-polar aliphatic groups retained or formed during pyrolysis \parencite{maoLinkingHydrophobicityBiochar2019}. Consequently, the apparent hydrophilicity of the soil decreased with increasing biochar content, as reflected by the progressive increase in contact angle. It should be noted that contact angle characterises the initial wetting response rather than the real water retention behaviour \parencite{zhangRolesBiocharsProperties2024}. Biochar surfaces are chemically heterogeneous, containing both hydrophobic carbonaceous domains and hydrophilic oxygen-containing functional groups. These polar functional groups facilitate water adsorption and may enhance soil water retention \parencite{zhangRolesBiocharsProperties2024, pardoMechanismImprovementBiochar2018}. With continued wetting, water can penetrate surface pores and interact with hydrophilic functional groups, leading to a gradual decrease in contact angle. This behaviour has recently been termed pseudo-hydrophobicity \parencite{jingDynamicContactAngle2026}.
\begin{figure}[htbp]
    \centering
    \includegraphics[width=1.0\linewidth]{../figures/Ch3/ch3CA.png}
    \caption{Contact angle images for (a) biochar, (b) untreated soil, (c) 5\% biochar-amended soil, and (d) 10\% biochar-amended soil}
    \label{fig:ch3CA}
\end{figure}

\subsection{Microstructural characteristics}
\subsubsection{\Ac{SEM} results}
Figure \ref{fig:ch3SEM} presents the SEM images of the tested specimens. According to the classification proposed by the Soil Science Society of America (SSSA), pores can be categorised as macropores (> 75 $\mu$m), mesopores (30–75 $\mu$m), micropores (5–30 $\mu$m), ultra-micropores (0.1–5 $\mu$m), and cryptopores (< 0.1 $\mu$m \parencite{dongMicrostructuralInvestigationUnsaturated2024}. For the untreated soil, soil particles tend to form aggregates, resulting in the development of substantial macropores and mesopores between adjacent aggregates, i.e., inter-aggregate pores. These pores are interconnected through pore throats, forming preferential pathways for water flow, as illustrated in Figure \ref{fig:ch3SEM}a. A higher-magnification image focusing on the aggregates (Figure \ref{fig:ch3SEM}b) reveals micropores, i.e., intra-aggregate pores, which in a range of around 5 $\mu$m. Some ultra-micropores among clay particles are also investigated. These micro- and ultra-micropores are crucial for capillary water retention. The clay particles form predominantly face-to-face and point-to-face contacts, giving rise to a metastable fabric that may undergo structural rearrangement under loading or suction changes.
For the biochar-amended soil, as shown in Figure \ref{fig:ch3SEM}c, sand, silt and biochar particles forms aggregated and interlocked structures, which can improve the aggregate stability \parencite{baiamonteEffectBiocharPhysical2019}. The sand particles are predominantly sub-rounded and in close contact with the biochar particles, while the interparticle voids are partially filled by silt particles and biochar fragments. At higher magnification (Figure \ref{fig:ch3SEM}d), mesopores and irregular macropores are identified at the biochar–soil interfaces, producing a heterogeneous pore network with tortuous flow paths. Plate-like biochar fragments are also observed. Figures \ref{fig:ch3SEM}e and \ref{fig:ch3SEM}f provide further evidence of clay-biochar interactions. On the one hand, clay particles are observed within biochar intrinsic pores, partially occluding the internal pore network \parencite{jingInteractionsBiocharClay2022}. Note that the diameter of these pores ranges from 1 to 5 $\mu$m. On the other hand, fractured intrinsic pores acts as attachment sites for clay particles due to the electrostatic attraction between the biochar surface and clay \parencite{yaoCharacterizationEnvironmentalApplications2014}. Due to clay's hydrophilic nature, its accumulation within and around the biochar pores can furtherincrease water retention potential. \par
\begin{figure}[htbp]
    \centering
    \includegraphics[width=1.0\linewidth]{../figures/Ch3/ch3SEM.png}
    \caption{SEM images of the tested materials: (a) untreated soil at $\times$ 500, (b) magnified view at $\times$ 2000, (c) biochar-amended soil at $\times$ 100, (d) magnified view at $\times$ 500, (e) magnified view at $\times$ 2000 highlighting soil particles retained within intact biochar intrinsic pores, and (f) attaching broken biochar intrinsic pores}
    \label{fig:ch3SEM}
\end{figure}
\subsubsection{\Ac{PSD} characteristics}
Figure \ref{fig:ch3PoreSize} presents the logarithmic differential \ac{PSD} ($\Delta V$/$\Delta\log D$) and cumulative mercury intrusion curves obtained from the \ac{MIP} test. Consistent with the \ac{SEM} observations, the pore-size is classified into macropores, mesopores, micropores, ultra-micropores and cryptopores, which are represented withdifferent background colours for clarity.\par
As depicted in Figures \ref{fig:ch3PoreSize}a and c, the compacted specimens exhibits bimodal \acp{PSD}, with a main pore peak ($D_{\mathrm{m}}$) in the ultra-micropore range (0.1 - 5 $\mu$m), and a secondary pore peak ($D_{\mathrm{s}}$) in the meso- and macropore ranges (> 75 $\mu$m).
For the densely compacted untreated soil (S1B0D), $D_{\mathrm{m}}$ is 0.43 $\mu$m, while $D_{\mathrm{s}}$ occurred between 10 and 20 $\mu$m. Based on the SEM observations, $D_{\mathrm{m}}$ is primarily associated with intra-aggregate pores, whereas $D_{\mathrm{s}}$ corresponds to inter-aggregate pores. Biochar has little influence on the \ac{PSD} within the cryptopore range, except for a slight increase in pore volume between 0.01 and 0.05 $\mu$m. With biochar, $D_{\mathrm{m}}$ within the ultra-micropore range shifts to 2.05 $\mu$m for S1B5D and 2.07 $\mu$m for S1B10D. For S1B5D, a weak residual peak remains at approximately 0.35 $\mu$m, close to that of the untreated soil. This feature disappears in S1B10D, where the \ac{PSD} increased continuously towards the dominant peak. Despite doubling the biochar content, $D_{\mathrm{m}}$ at approximately 2 $\mu$m remained unchanged in both position and intensity, together with almost overlapping descending branches of the \ac{PSD} curves. This suggests that the 2 $\mu$m peak is unlikely to caused by the intrinsic pores of the biochar. Instead, it is associated with stable structural pores formed by the interaction between biochar and soil particles. The micropore and mesopore ranges (approximately 6–50 $\mu$m) are highly sensitive to biochar content. At a biochar content of 5\%, the pore volume in this range decreased to approximately one-third of that of the untreated soil. Increasing the biochar content to 10\% resulted in a partial recovery of pore volume within this size range. These observations indicate that rather than generating additional macropores, biochar appears to reorganise the original 10–20 $\mu$ pore system into a stable aggregate-associated pore domain of approximately 0.3–2 $\mu$m. The delimiting pore size ($D_{\mathrm{L}}$) of all three specimens is approximately 5 $\mu$m, corresponding to the boundary between the ultra-micropore and micropore ranges.\par
Compared with the densely compacted specimens, the \ac{PSD} of the loosely compacted specimens exhibits a more pronounced bimodal pattern (Figure \ref{fig:ch3PoreSize}c). For the untreated soil (S1B0L), $D_{\mathrm{m}}$ is centred at 0.82 $\mu$m. With biochar, $D_{\mathrm{m}}$ slightly increases to 2.15 $\mu$m and 2.13 $\mu$m for S1B5L and S1B10L, respectively. The $D_{\mathrm{L}}$ shows slight increase for all loosely compacted specimens. These results suggest that the pore system of soil–biochar aggregate is relatively stable, which, together with the SEM observations, can be attributed to the angular morphology of the biochar particles that enhances particle interlocking. In contrast, the larger pore domain is much more affected by biochar dosage. When the biochar content increases to 5\% and 10\%, $D_{\mathrm{s}}$ decreases from 60.05 $\mu$m for S1B0L to 29.84 $\mu$m for S1B5L and 22.5 $\mu$m for S1B10L. However, an increase in pore volume is also observed in the macropore region, particularly for pores larger than 156 $\mu$m. Figure \ref{fig:ch3porecategory} presents the proportions of different pore categories within the detected pore size range. At 90\% \ac{DC}, the addition of biochar introduces more ultra-micropores and less micropores, while the proportion of other pore categories shows slight changes. However, at 80\% \ac{DC}, only micropores shows a slight slight increase with biochar dosage, while the proportions of other pore categories remain relatively stable. This indicates that biochar mainly affects the intra-aggregate pores, which includes ultra-micropores and micropores, while the inter-aggregate pores are less affected. With increasing \ac{DC}, the mesopore and macropore domains were affected most significantly, accounting for the majority of the reduction in total pore volume. As demonstrated in Figure \ref{fig:ch3SEM}c and d, biochar shows heterogeneous role in the formation of the soil–biochar pore system. On the one hand, smaller biochar fragments fill inter-aggregate voids, leading to a shift of ($D_{\mathrm{s}}$) from the mesopore range towards the micropore range. On the other hand, those larger biochar particles, which not fully embedded within soil aggregates, introduce additional macropores due to their irregular geometry.\par
Figures \ref{fig:ch3PoreSize}b and d present the cumulative mercury intrusion curves obtained from the \ac{MIP} tests. As expected, reducing the \ac{DC} results in a higher cumulative intrusion volume, indicating a greater volume of mercury-accessible pores. At a given \ac{DC}, the cumulative intrusion volume also increases with biochar content, suggesting that biochar incorporation increases the accessible pore volume. Table \ref{tab:ch3_pore_structure_parameters} compares the total void ratio determined during specimen preparation with that calculated from the MIP measurements. The MIP-derived void ratio is consistently 3–12\% lower. This discrepancy is commonly attributed to the presence of isolated or poorly connected pores, as well as the ink-bottle effect, whereby narrow pore throats restrict mercury access to larger pore bodies. With the increase of biochar content, the intruded void ration increases, indicating that biochar likely enhances pore connectivity. Furthermore, the densely compacted specimens exhibit a larger proportion of mercury-inaccessible pores than the loosely compacted specimens, likely because compaction reduces pore connectivity and constricts pore throats, preventing mercury from accessing part of the pore network.\par
\begin{figure}[htbp]
    \centering
    \includegraphics[width=1.0\linewidth]{../figures/Ch3/ch3PoreSize.png}
    \caption{Measured pore size distribution from \ac{MIP} test of densely compacted speciments ((a) and (b)), and loosely compacted specimens ((c) and (d))}
    \label{fig:ch3PoreSize}
\end{figure}


\begin{figure}[htbp]
    \centering
    \includegraphics[width=1.0\linewidth]{../figures/Ch3/ch3porecategory.png}
    \caption{Percentage of different pore categories in the tested specimens}
    \label{fig:ch3porecategory}
\end{figure}




\begin{table}[htbp]
    \centering
    \caption{Pore structure parameters of representative specimens}
    \label{tab:ch3_pore_structure_parameters}
    \begingroup
    \renewcommand{\arraystretch}{1.15}
    \setlength{\tabcolsep}{2pt}
    \begin{tabular*}{\textwidth}{@{\extracolsep{\fill}}>{\centering\arraybackslash}m{0.12\textwidth}>{\centering\arraybackslash}m{0.11\textwidth}>{\centering\arraybackslash}m{0.15\textwidth}>{\centering\arraybackslash}m{0.18\textwidth}>{\centering\arraybackslash}m{0.10\textwidth}>{\centering\arraybackslash}m{0.10\textwidth}>{\centering\arraybackslash}m{0.10\textwidth}@{}}
        \toprule
        \makecell[c]{Specimen\\ID} & \makecell[c]{Void\\ratio} & \makecell[c]{Intruded\\void ratio} & \makecell[c]{Non-detected\\void ratio} & \makecell[c]{$D_{\mathrm{m}}$\\($\mu$m)} & \makecell[c]{$D_{\mathrm{L}}$\\($\mu$m)} & \makecell[c]{$D_{\mathrm{s}}$\\($\mu$m)} \\
        \midrule
        S1B0D  & 0.599 & 0.532 & 0.067 & 0.43 & 5.33 & 18.87 \\ 
        S1B0L  & 0.802 & 0.729 & 0.073 & 0.82  & 6.52 & 60.05 \\ 
        S1B5D  & 0.59  & 0.548 & 0.042 & 2.05  & 5.17 & 44.91 \\ 
        S1B5L  & 0.803 & 0.767 & 0.036 & 2.15  & 7.17 & 29.84 \\ 
        S1B10D & 0.603 & 0.544 & 0.059 & 2.07  & 4.82 & 30.62 \\ 
        S1B10L & 0.807 & 0.78  & 0.026 & 2.13  & 6.52 & 23.27 \\ 
        \bottomrule
    \end{tabular*}
    \endgroup
\end{table}
\subsubsection{Fractal dimension analysis}
For a three dimensional pore system, the fractal dimension of the pore surface ($d_{\mathrm{s}}$) is widely used to characterize the self-similarity and complexity of pore systems \parencite{perfectApplicationsFractalsSoil1995, hanRelationshipFractalFeature2022}. In this study, $d_{\mathrm{s}}$ is dertermined by the thermodynamic fractal model, which has been demonstrated to be more suitable for characterising the pore structure of soils \parencite{shiStrengthMicroscopicPore2024}. In \ac{MIP} test, the applied pressure overcomes the capillary resistance that prevents the mercury from entering pore throats. The increase in surface free energy of mercury originates from the external work performed during mercury intrusion. By assuming pore system as an equivalent bundle of ideal capillary tubes, the relationship between the external pressure and the incremental volume of intruded mercury can be expressed as:
\begin{equation}
     \int_{0}^{V_{\mathrm{m}}}p_{\mathrm{m}}\mathrm{d}V=-\int_{0}^{S_{\mathrm{m}}}\gamma \cos\theta \mathrm{d}S
\end{equation}
which can be further reformed as \parencite{shiStrengthMicroscopicPore2024}:
\begin{equation}
    \label{eq:ch3mip2}
    \sum_{i=1}^{n}p_{\mathrm{m},n}\Delta V_{n}=kr_{n}^{2}\left ( V_{n}^{1/3}/r_{n} \right )
\end{equation}
where $S$ is the surface area of the mercury intrusion pores, $p_{\mathrm{m},n}$ is the applied pressure at the $n$-th intrusion step, $V_{n}$ and $\Delta V_{n}$ represents the accumulative and incremental volume of mercury at the $n$-th intrusion step, $k$ is a constant, and $r_{n}$ is the equivalent radius of the pore. Note that the term in the left hand side of the Equation \ref{eq:ch3mip2} represents the total surface free energy of the mecury, while $V_{n}^{1/3}/r_{n}$ represents change of mecury volume. Taking the logarithm of both sides, Equation \ref{eq:ch3mip2} can be further expressed as:
\begin{equation}
    \log\left[\left(\sum_{i=1}^{n}p_{\mathrm{m},i}\Delta V_{i} \right)/r_{n}^{2} \right] = d_{\mathrm{s}} \log \left( V_{n}^{1/3}/r_{n} \right) + R
\end{equation}
where $d_{\mathrm{s}}$ and $R$ can be determined as the slope and intercept of the regression line, respectively, as presented in Figure \ref{fig:ch3FractalLine}. The results of fractal dimension are summarised in Figure \ref{fig:ch3fractalDimensionValue}. At all biochar contents, $d_{\mathrm{s}}$ is higher at 90\% than at 80\% \ac{DC}. As investigated in \ac{PSD} results, compaction compresses large pores, promotes particle rearrangement, and enhances soil–biochar interfacial interactions, causing the remaining pores to evolve towards a finer, more irregular, and heterogeneous multiscale structure \parencite{sunFractalCharacteristicsPore2020, mengWaterholdingCharacteristicsLoess2026}. At 90\% \ac{DC}, $d_{\mathrm{s}}$ initially increased and then decreased slightly with increasing biochar content. This trend was consistent with the variation in the amplitude of the \ac{PSD} curve near the secondary peak $D_{\mathrm{s}}$ (see Figure \ref{fig:ch3PoreSize}a). As further shown in Figure \ref{fig:ch3porecategory}, S1B5D exhibited the greatest heterogeneity in the distribution of pore categories. In contrast, at 80\% \ac{DC}, the separation between $D_{\mathrm{m}}$ and $D_{\mathrm{s}}$ progressively decreased with increasing biochar content, indicating a more uniform pore-size distribution and, correspondingly, a reduction in $d_{\mathrm{s}}$.\par
\begin{figure}[htbp]
    \centering
    \includegraphics[width=1.0\linewidth]{../figures/Ch3/ch3FractalLine.png}
    \caption{Regression analysis of fractal dimension of the tested speciments: (a) 90\% \ac{DC}, and (b) 80\% \ac{DC}}
    \label{fig:ch3FractalLine}
\end{figure}

\begin{figure}[htbp]
    \centering
    \includegraphics[width=0.6\linewidth]{../figures/Ch3/ch3fractalDimensionValue.png}
    \caption{Fractal dimension values of the tested specimens}
    \label{fig:ch3fractalDimensionValue}
\end{figure}

\subsection{\Ac{SWRC} at zero stress}
Figures \ref{fig:ch3HYPROPWP4CSWRC}a and b present the drying \ac{SWRC}s of the untreated and biochar-amended soils compacted to 90\% and 80\% \ac{DC}, respectively. The curves are obtained by combining HYPROP measurements in the low-suction range ($\leq$ 100 kPa), with WP4C measurements in the high-suction range (> 100 kPa). The experimental data are fitted using the \ac{VG} model \parencite{vangenuchtenClosedformEquationPredicting1980}, which can be expressed as:
\begin{equation}
    \theta = \theta_{\mathrm{r}}+\frac{\theta_{\mathrm{s}} - \theta_{\mathrm{r}}}{[1+(\alpha \cdot s)^{n}]^{m}}
\end{equation}
where $\theta _{\mathrm{s}}$ is the saturated volumetric water content, $\theta _{\mathrm{r}}$ is the residual volumetric water content, $\alpha$, $m$, and $n$ are the fitting parameters. The Casagrande method was used to determine the \ac{AEV} \parencite{vanapalliInfluenceSoilStructure1999}. The fitting parameters and air-entry values are summarised in Table \ref{tab:ch3_vgfitting}.\par
For both densely and loosely compacted specimens, the addition of biochar causes only minor changes in the \ac{SWRC} under near-saturated conditions, particularly at matric suctions below 10 kPa. With the increase in matric suction, the biochar-amended specimens exhibit enhanced water retention capacity. A comparison between Figures \ref{fig:ch3HYPROPWP4CSWRC}a and b further shows that biochar has a more pronounced effect on the \ac{SWRC} under dense compaction than under loose compaction. This behaviour can be interpreted from the \ac{SEM} and \ac{PSD} results. The densely compacted soil initially contains a lower proportion of inter-aggregate mesopores and macropores. In contrast, the loosely compacted soil already contains a relatively developed inter-aggregate pore network. Therefore, increasing the biochar content leads to a more pronounced modification of pore structure, making the corresponding \ac{SWRC} more sensitive to biochar dosage.\par
The \ac{AEV} represents a critical matric suction at which air begins to enter the pore network, marking the onset of significant desaturation. At 90\% \ac{DC}, the \ac{AEV} increases from 4.986 kPa for the untreated soil to 6.152 kPa with 5\% biochar addition. However, a further increase in biochar content to 10\% reduces the \ac{AEV} to 5.017 kPa. A different trend is observed for the loosely compacted specimens. At 80\% \ac{DC}, biochar incorporation consistently increases the \ac{AEV} from 3.017 kPa for S1B0L to 3.894 kPa for S1B5L and 4.258 kPa for S1B10L. According to the Young–Laplace equation, the equivalent pore size corresponding to the 
\ac{AEV} can be calculated as:
\begin{equation}
    D_{\mathrm{AEV}} = \frac{4 \gamma \cdot \cos \theta}{\mathrm{AEV}}
\end{equation}
where $\gamma$ is the surface tension of water, $\theta$ is the contact angle (see Section \ref{sec:ch3contactangle}), and $\mathrm{AEV}$ is the air-entry value. The results of $D_{\mathrm{AEV}}$ are also listed in Table \ref{tab:ch3_vgfitting}. As listed in Table \ref{tab:ch3_vgfitting}, $D_\mathrm{AEV}$ of the specimens compacted to 90\% are 57.34, 44.62, and 53.25 $\mu$m for S1B0D, S1B5D, and S1B10D, respectively. The variation in \ac{AEV} can therefore be related to changes in the pore-throat population around this range. As shown in Figure \ref{fig:ch3PoreSize}a, within the pore size less than 50 $\mu$m, the \ac{PSD} intensity follows the order S1B0D > S1B10D > S1B5D. The addition of 5\% biochar reduces the abundance of pore throats in this range, leading to a higher \ac{AEV}, whereas increasing the biochar content to 10\% partially restores this pore-throat frequency and consequently lowers the \ac{AEV}. A similar relationship is also observed for the specimens compacted to 80\%. According to Table \ref{tab:ch3_vgfitting}, the variation in \ac{AEV} can be interpreted in terms of the pore-throat approximately 40-100 $\mu$m. As shown in Figure \ref{fig:ch3PoreSize}c, biochar incorporation progressively reduces the abundance of pore throats within this size range. Correspondingly, the \ac{AEV} increases with increasing biochar content.\par
In the \ac{VG} model, the fitting parameter $n$ describes the slope of \ac{SWRC} during drying. As summarised in Table \ref{tab:ch3_vgfitting}, increasing the biochar content decreases the $n$ value at both compaction degrees, indicating a more gradual reduction in water content with increasing matric suction and, consequently, enhanced water retention over a wider suction range. A similar trend has been reported in previous studies and is generally attributed to biochar-induced pore refinement \parencite{dongMicrostructuralInvestigationUnsaturated2024,hussainInfluenceBiocharSoil2020, hussainInvestigatingUnsaturatedHydraulic2021}. The \ac{PSD} results in the present study also support this interpretation. Specifically, biochar increases the abundance of micropores while shifting the characteristic mesopore and macropore domains towards smaller pore sizes. This refinement of the pore system broadens the distribution of water-filled pores and delays pore drainage during drying, resulting in a less steep \ac{SWRC}.\par
\begin{figure}[htbp]
    \centering
    \includegraphics[width=1.0\linewidth]{../figures/Ch3/ch3HYPROPWP4CSWRC.png}
    \caption{\ac{SWRC} in the evaporation tests: (a) densely compacted specimens, and (b) loosely compacted specimens}
    \label{fig:ch3HYPROPWP4CSWRC}
\end{figure}

\begin{table}[htbp]
    \centering
    \caption{Fitting parameters of the \ac{VG} model and \ac{AEV}}
    \label{tab:ch3_vgfitting}
    \begingroup
    \small
    \renewcommand{\arraystretch}{1.2}
    \setlength{\tabcolsep}{1pt}
    \begin{tabularx}{0.97\textwidth}{@{}>{\centering\arraybackslash}m{0.105\textwidth}*{7}{>{\centering\arraybackslash}X}>{\centering\arraybackslash}m{0.15\textwidth}@{}}
        \toprule
        \makecell[c]{Specimen\\ID} & $\theta_{\mathrm{r}}$ & $\theta_{\mathrm{s}}$ & \makecell[c]{$\alpha$\\(kPa$^{-1}$)} & $n$ & $m$ & $R^2$ & \makecell[c]{AEV\\(kPa)} & \makecell[c]{Equivalent\\pore-throat\\diameter\\($\mu$m)} \\
        \midrule
        S1B0D  & 0.0853 & 0.374 & 0.105 & 1.128 & 0.113 & 0.994 & 4.986 & 57.34 \\
        S1B0L  & 0.0852 & 0.445 & 0.144 & 1.240 & 0.194 & 0.998 & 3.017 & 93 \\
        S1B5D  & 0.1160 & 0.368 & 0.085 & 1.125 & 0.111 & 0.995 & 6.152 & 44.62 \\
        S1B5L  & 0.0851 & 0.445 & 0.136 & 1.175 & 0.149 & 0.997 & 3.894 & 70.47 \\
        S1B10D & 0.0850 & 0.369 & 0.104 & 1.077 & 0.072 & 0.992 & 5.017 & 53.25 \\
        S1B10L & 0.0852 & 0.447 & 0.121 & 1.162 & 0.139 & 0.995 & 4.285 & 62.35 \\
        \bottomrule
    \end{tabularx}
    \endgroup
\end{table}

\subsection{Hydraulic conductivity}
Figure \ref{fig:ch3HC} presents the relationship between the hydraulic conductivity ($k_{\mathrm{un}}$) and matric suction measured in the first series of tests. At the begining of the test, the specimens are fully saturated, with all pores occupied by free water. As drying progressed, $k_{\mathrm{un}}$ decreases rapidly, particularly after the matric suction exceeds the \ac{AEV} (3 to 10 kPa). This is because air first invades the larger pores, reducing effective cross-sectional area available for water flow and increasing the tortuosity of the remaining flow paths \parencite{vangenuchtenClosedformEquationPredicting1980,fredlundPredictingPermeabilityFunction1994}. With continued drying, water is prone to forming thin films coating the soil and biochar particles or retained within smaller pores. Continued drying reduced thickness and hydraulic conductance of water films, thereby increasing viscous resistance and decreasing $k_{\mathrm{un}}$ \parencite{tullerHydraulicConductivityVariably2001}.\par
At matric suctions below 10 kPa, the change of $k_{\mathrm{un}}$ can be interpreted by \ac{PSD} and cumulative pore volume, as demonstrated in Figure \ref{fig:ch3PoreSize} and Table \ref{tab:ch3_pore_structure_parameters}. The 80\% \ac{DC} specimens exhibits higher $k_{\mathrm{un}}$ than the 90\% \ac{DC} specimens because of their larger fractions of macropores and mesopores. S1B0L, with the largest secondary pore peak ($D_{\mathrm{s}}$ = 60.05 $\mu\mathrm{m}$) and the highest cumulative macropore–mesopore fraction (approximately 30\%), presents the highest $k_{\mathrm{un}}$, but also the most rapid decline when the matric suction exceeded the \ac{AEV}. For the 80\% \ac{DC} specimens, biochar reduces $D_{\mathrm{s}}$ from 60.05 $\mu\mathrm{m}$ to 29.84 and 23.27 $\mu\mathrm{m}$ and lowers the fractions of mesopores and macropores, indicating refinement of the principal flow paths. Accordingly, $k_{\mathrm{un}}$ decreases in the order S1B0L > S1B5L > S1B10L, consistent with the pattern of\ac{PSD}. Among the 90\% \ac{DC} specimens, S1B5D has the lowest pore volume over the range 6–50 $\mu\mathrm{m}$. And its higher fractal dimension ($d_{s}$ = 2.7943) indicates a more tortuous pore network. The macro-mesopore fraction is highest in S1B0D, followed by S1B10D and S1B5D. The corresponding $k_{\mathrm{un}}$ follows the same order, i.e., S1B0D > S1B10D > S1B5D.\par
At matric suctions above 10 kPa, all specimens exceed their respective \acp{AEV}. As water drains from the macropores and mesopores, $k_{\mathrm{un}}$ becomes increasingly controlled by the finer pore network. Under both \acp{DC}, the dominant pore diameters of the biochar-amended specimens are nearly identical (2.05 and 2.07 $\mu\mathrm{m}$ for the 90\% \ac{DC} specimens and 2.15 and 2.13 $\mu\mathrm{m}$ for the 80\% \ac{DC} specimens). Accordingly, $k_{\mathrm{un}}$ curves of B5 and B10 almost coincide beyond 40 kPa. Under both compaction conditions, B10 has a lower fractal dimension, indicating less tortuous and better-connected residual flow paths. Consequently, B10 exhibits a slightly higher $k_{\mathrm{un}}$ than B5 at suctions above 40 kPa. In contrast, although the untreated soil exhibits the highest $k_{\mathrm{un}}$ at low suctions, its poorer water-retention capacity causes the fine-scale water pathways to lose continuity more readily during drying, resulting in the lowest $k_{\mathrm{un}}$ at suctions above 40 kPa.\par
\begin{figure}[htbp]
    \centering
    \includegraphics[width=1.0\linewidth]{../figures/Ch3/ch3HC.png}
    \caption{Hydraulic conductivity of the tested specimens: (a) 90\% \ac{DC}, and (b) 80\% \ac{DC}}
    \label{fig:ch3HC}
\end{figure}

\subsection{\Ac{SD-SWRC} and degree of hysteresis}




\subsection{One-dimentional compressibility}


\section{Summary}
summary check